% =====================================================================
% v2 paper main TeX file — IEEE conference style skeleton
% =====================================================================
%
% This file is the assembly point. The 7 paper section files are written
% as Markdown under cloud_removal_v2/docs/paper_v2_section_*.md and must
% be converted to LaTeX before \input. Two recipes are documented in
% Build instructions below; pick whichever your environment supports.
%
% Build instructions
% ------------------
%
% Recipe A — pandoc (recommended, single command):
%
%   for s in 1_intro 2_related_work 3_system_model 4_theory \
%            5_system_opt 6_A_setup 6_B_noniid 6_C_fedbn \
%            6_D_aggregation 6_E_ann_vs_snn 6_F_qualitative \
%            6_G_per_plane 6_H_discussion 7_conclusion; do
%     pandoc paper_v2_section_$s.md \
%            --from markdown --to latex \
%            --output paper_v2_section_$s.tex
%   done
%
% Recipe B — manual conversion (no pandoc):
%   For each section, copy the Markdown content into a .tex file and
%   apply: ## -> \subsection, ** ** -> \textbf{}, * * -> \emph{},
%   pipe-tables -> tabular, math $..$ stays the same.
%
% Then build:
%   pdflatex paper_v2_main.tex
%   bibtex   paper_v2_main
%   pdflatex paper_v2_main.tex
%   pdflatex paper_v2_main.tex
%
% =====================================================================

\documentclass[conference]{IEEEtran}
\IEEEoverridecommandlockouts

% --- Core packages ---
\usepackage[utf8]{inputenc}
\usepackage[T1]{fontenc}
\usepackage{cite}
\usepackage{amsmath,amssymb,amsfonts,amsthm}
\usepackage{algorithmic}
\usepackage{algorithm}
\usepackage{graphicx}
\usepackage{textcomp}
\usepackage{xcolor}
\usepackage{booktabs}
\usepackage{multirow}
\usepackage{hyperref}
\usepackage{url}
\usepackage{listings}

% --- Hyperref colour conventions ---
\hypersetup{
  colorlinks   = true,
  linkcolor    = black,
  citecolor    = blue,
  urlcolor     = blue,
}

% --- Theorem environments (used in section IV) ---
\newtheorem{theorem}{Theorem}
\newtheorem{proposition}{Proposition}
\newtheorem{corollary}{Corollary}
\newtheorem{lemma}{Lemma}
\newtheorem{assumption}{Assumption}
\newtheorem{claim}{Claim}
\newtheorem{remark}{Remark}

% --- Useful macros ---
\newcommand{\Var}{\mathrm{Var}}
\newcommand{\E}{\mathbb{E}}
\newcommand{\R}{\mathbb{R}}
\newcommand{\Dir}{\mathrm{Dir}}
\newcommand{\Beta}{\mathrm{Beta}}
\newcommand{\D}{\mathcal{D}}
\newcommand{\N}{\mathcal{N}}
\newcommand{\Iset}{\mathcal{I}}
\newcommand{\eff}[1]{\mathrm{eff}(#1)}

% --- Path for figures (PDFs go under figures/) ---
\graphicspath{{figures/}}

% =====================================================================
\begin{document}

% --- Title block ---
\title{Brain-Inspired Decentralised Satellite Federated Learning for
       Cloud Removal under Source-Level Dirichlet Non-IID:
       FedBN-Redundancy and the Boundary Case of FLSNN's Theorem~2}

% Author block — replace with actual authors and affiliations.
\author{
  \IEEEauthorblockN{(authors anonymised for review)}\\
  \IEEEauthorblockA{(institutions anonymised for review)}
}

\maketitle

% =====================================================================
% Abstract — content from cloud_removal_v2/docs/paper_v2_abstract.md
% =====================================================================
\begin{abstract}
We extend the FLSNN brain-inspired decentralised satellite
federated-learning framework from on-board image classification
to the more demanding setting of pixel-level cloud removal under
source-level Dirichlet non-IID. Inheriting the 50/5/1
Walker-Star constellation, intra-plane ring-AllReduce, and the
inter-plane RelaySum / Gossip / AllReduce comparison verbatim,
we contribute two theoretical and three empirical results.
Theoretically, we derive a closed-form upper bound on the
inter-plane gradient dissimilarity $\zeta^2$ for an $S$-source
Dirichlet partition (Proposition~\ref{prop:dirichlet}) and show
that, in the resulting small-$\zeta^2$ regime, the three
aggregation schemes' Theorem-\ref{thm:flsnn} bounds collapse to
the same asymptotic order (Corollary~\ref{cor:collapse}).
Empirically, on the CUHK-CR benchmark with $\alpha = 0.1$,
AllReduce Pareto-dominates Gossip and RelaySum at $60\,\%$ less
inter-plane communication, an ANN backbone beats SNN by $+0.75$
\,dB PSNR / $1.61\times$ wall time on a single GPU, and
FedBN-style BN-local aggregation yields gains within the
single-seed noise floor regardless of BN variant --- TDBN
reduces cross-plane $\Var(\gamma)$ by $42\,\%$ vs standard BN,
but neither variant produces drift large enough for FedBN to
repair. The five findings are mutually consistent with FLSNN's
Theorem~\ref{thm:flsnn} in the small-$\zeta^2$ boundary regime.
\end{abstract}

\begin{IEEEkeywords}
Spiking neural networks, satellite federated learning, cloud
removal, Dirichlet non-IID, batch normalisation, decentralised
optimisation.
\end{IEEEkeywords}

% =====================================================================
% Section content — \input the converted .tex files. If pandoc is not
% used, replace each \input with the inline section content.
% =====================================================================

\section{Introduction}\label{sec:intro}
\input{paper_v2_section_1_intro.tex}

\section{Related Work}\label{sec:related}
\input{paper_v2_section_2_related_work.tex}

\section{System Model and Algorithm}\label{sec:sysmodel}
\input{paper_v2_section_3_system_model.tex}

\section{Theoretical Analysis}\label{sec:theory}
\input{paper_v2_section_4_theory.tex}

\section{Inter-Plane Topology Optimisation}\label{sec:sysopt}
\input{paper_v2_section_5_system_opt.tex}

\section{Experiments}\label{sec:experiments}
\input{paper_v2_section_6_A_setup.tex}
\input{paper_v2_section_6_B_noniid.tex}
\input{paper_v2_section_6_C_fedbn.tex}
\input{paper_v2_section_6_D_aggregation.tex}
\input{paper_v2_section_6_E_ann_vs_snn.tex}
\input{paper_v2_section_6_F_qualitative.tex}
\input{paper_v2_section_6_G_per_plane.tex}
\input{paper_v2_section_6_H_discussion.tex}

\section{Conclusion}\label{sec:conclusion}
\input{paper_v2_section_7_conclusion.tex}

% =====================================================================
% Bibliography
% =====================================================================
\bibliographystyle{IEEEtran}
\bibliography{paper_v2_refs}

% =====================================================================
% Appendix (optional) — full proof of Proposition 1
% =====================================================================
\appendix
\section*{Appendix A: Notation}
See Table~III.F (notation summary, in \texttt{paper\_v2\_section\_3\_system\_model.md}).

\section*{Appendix B: Proof of Proposition~\ref{prop:dirichlet}}
The full derivation chain (Setup, Gradient decomposition,
Beta($\alpha,\alpha$) moments, expected average squared
dissimilarity, sup-norm form) is in
\texttt{cloud\_removal\_v2/docs/v2\_formal\_derivations.md}
\S\,D1.1--D1.6. We provide the headline computation here for
the camera-ready draft.

\medskip
\noindent\textbf{Setup recap.} Each client's per-sample distribution
is $p_i \D_1 + (1-p_i) \D_2$ with $p_i \sim \Beta(\alpha,\alpha)$ i.i.d.
across $i = 1, \ldots, N$.

\medskip
\noindent\textbf{Gradient decomposition.} By linearity of
$\nabla_\theta$ in $p_i$:
$\nabla f_i - \nabla f = (p_i - \bar p_N)(\nabla f_1 - \nabla f_2)$.
Squaring:
$\| \nabla f_i - \nabla f \|^2
= (p_i - \bar p_N)^2 \|\nabla f_1 - \nabla f_2\|^2$.

\medskip
\noindent\textbf{Beta moments.} $\E[p_i] = 1/2$;
$\Var_{\Beta(\alpha,\alpha)}(p_i) = 1/(4(2\alpha+1))$.

\medskip
\noindent\textbf{Expected average.} Using
$\E[\sum_i (p_i - \bar p_N)^2] = (N-1)\Var(p_i)$:
$$
\E_{\mathbf{p}}\!\Bigl[\tfrac{1}{N}\sum_i (p_i - \bar p_N)^2\Bigr]
= \frac{N-1}{4N(2\alpha+1)}.
$$
Combining with the gradient decomposition gives
Proposition~\ref{prop:dirichlet}. $\hfill\square$

\section*{Appendix C: Build instructions}

This file's compilation requires \texttt{IEEEtran.cls} (IEEE
conference style) and a \LaTeX{} distribution with the standard
\texttt{amsmath, booktabs, multirow, hyperref, graphicx} packages.

The 14 referenced figures are produced from the saved checkpoints
and \texttt{summary.json} files via the seven plotting scripts
listed in \texttt{cloud\_removal\_v2/docs/tables/README.md}
\S\,7. Place them under \texttt{figures/} and update each
section's \texttt{\textbackslash includegraphics\{...\}} path.

The full reproducibility chain (every numeric cell of the paper's
4 tables) is in \texttt{paper\_v2\_MASTER\_INDEX.md} \S\,7
(``End-to-end reproduction recipe'').

\end{document}
